\section{Introduction}
\label{sec:ch4-introduction}
Interactive systems are integral to building design, leveraging advances in sensing, actuation, and data processing to regulate parameters such as lighting, temperature, and air quality~\cite{alavi2017comfort}. These systems enhance comfort, energy efficiency, and sustainability while adapting spatial layouts to create private or collaborative spaces within shared environments~\cite{nguyen2022towards, nguyen2024adaptive, deng2019diffusive}. However, this approach often embodies an anthropocentric and technocentric paradigm~\cite{gill2008rethinking}, treating buildings---and the technologies within them---as functional spaces optimised for productivity~\cite{heidegger2006building}. This perspective neglects the potential to enrich deeper, experiential dimensions, such as those rooted in our innate connection to nature~\cite{wilson1986biophilia, kellert1995biophilia, kellert2018nature}.


Biophilia refers to the intrinsic human tendency to seek connections with nature and other forms of life~\cite{wilson1986biophilia}. Biophilic design builds on this, advocating for the integration of natural elements, patterns, and processes into the built environment to support meaningful connections with nature. This approach is grounded in a growing body of evidence demonstrating that a biophilic connection yields psychological and physiological benefits, including stress reduction, cognitive restoration, and enhanced creativity~\cite{frumkin2017nature, hartig2014nature}.


In architectural legacy, biophilia is deeply embedded. Traditional methods have employed natural materials and site-sensitive layouts to promote ecological awareness, sensory engagement, and emotional connection~\cite{browning2020nature, kellert2018nature}. For instance, an ancient Indian temple at a river’s origin cyclically submerges during monsoon season, symbolising purification and renewal---of both the building and the soul---while reinforcing human-ecological ties~\cite{geva2023water}. Similarly, Frank Lloyd Wright’s Fallingwater in Pennsylvania exemplifies this, blending the built form with the surrounding waterfall\sidenote{https://fallingwater.org/}.

Despite this legacy, biophilic design remains underutilised in smart building technologies, which rarely explore how interactive systems might foster deeper, multisensory relationships with nature. There is a need to integrate biophilic principles into smart building research, towards environments that are not only advanced but also ecologically and experientially rich. Human-Building Interaction (HBI), which explores technologies, spaces, and subjective experiences~\cite{alavi2019introduction, becerik2022field, rao2023towards}, offers a promising context for this integration. By acknowledging emotional, cultural, and sensory dimensions, HBI can support environments that promote well-being, nature-connectedness, and sustainability for ``human flourishing''~\cite{mccormack2014no}.

This paper explores the following question: \textit{What insights can be drawn from the legacy, practice, and imaginaries of biophilic architecture to guide the design of interactive experiences in future buildings?} Through interviews with 13 architecture experts, we investigate the potential of biophilic design within the context of current and future smart buildings. We acknowledge biophilic design as a reciprocal practice benefiting both humans and nature, though current discourse remains largely human-centred; our perspective focuses on expert insights rather than normative prescriptions.

Our analysis reveals a gap in current practices and the possibilities of nature-connected interactions in buildings. We present \emph{seven high-level topics (TO1–TO7)} that offer a structured overview of expert discourse, followed by \emph{eight themes (T1–T8)} that capture how professionals imagine the future of biophilic design. These themes provide a conceptual foundation for nature-inspired interactions in indoor spaces. Next, we introduce \emph{five design opportunities (DO1–DO5)} supported by emerging technologies to enable novel biophilic interactions. By bridging architectural biophilia with interactive technologies, we expand the discourse on smart buildings---reimagining how nature can be integrated into indoor experiences and reorienting biophilic interaction as reciprocal and aligned with more-than-human, justice-oriented futures.

